\documentclass[11pt,a4paper]{article}
\usepackage{jambronotes}
\renewcommand{\memouser}{ACB}      % default username (your initials)
\renewcommand{\memocolor}{BoEred}  % default color

% --- Demo-only helper (NOT part of the theme) -------------------------------
% Shows the LaTeX source of each feature directly above its rendered result so
% readers can copy the exact usage. Set on the same light-gray background as the
% theme's `code` environment (codegray + jambrocode style) so source and result
% share a consistent look. Defined here, not in the .sty (demo-only).
\newtcblisting{democode}{%
  listing only,
  listing options={style=jambrocode},
  colback=codegray, colframe=codegray,   % grey fill, no visible border
  boxrule=0pt, arc=2pt,
  left=4pt, right=4pt, top=3pt, bottom=3pt,
  before skip=6pt, after skip=6pt,
  breakable,
}

%==================================================
\title{\vspace{-2cm}\color{title}
\huge\textbf{jambronotes.sty}\\[0.3em]
\Large Feature Demo\\[0.3em]
\large v1.1-dev --- \textit{\today}
}
\author{}
\date{}
\begin{document}
\maketitle
\tableofcontents
\clearpage
\listofboxes
\clearpage
%==================================================

%==================================================
\section{Typography and Fonts}
%==================================================

\subsection{Body Text}
%--------------------------------------------------

Body text is set in \textbf{Source Sans Pro} (loaded via \verb|[default]{sourcesanspro}| as the document default), while section headings use \textbf{Fira Sans} (the \verb|\firatitle| switch) for contrast. Captions follow the body font. Math is rendered in the standard Computer Modern math fonts.

A short paragraph: The quick brown fox jumps over the lazy dog. Numbers: 0123456789. Greek: $\alpha, \beta, \gamma, \delta, \epsilon$. Operators: $\sum_{i=1}^{n} x_i$, $\int_0^\infty e^{-x}\,dx$.

\subsection{Handwriting Font}
%--------------------------------------------------

The Augie font is available for hand-drawn effects. \texttt{\string\hand} is a font switch (use inside a group \verb|{...}|); \texttt{\string\handbold} takes its text as an argument:

\begin{democode}
{\hand This is the Augie handwriting font.}
\handbold{This is handbold (default stroke 0.6).}
\handbold[0.3]{Lighter stroke.}   % optional width argument
\handbold[1.0]{Heavier stroke.}
\end{democode}

\noindent renders as:

\begin{itemize}
    \item {\hand This sentence is set in the Augie handwriting font.}
    \item \handbold{This sentence uses handbold — bold Augie via PDF stroke literal.}
    \item Stroke width is an optional argument: \handbold[0.2]{light (0.2)}, \handbold[0.6]{default (0.6)}, \handbold[1.1]{heavy (1.1)}.
    \item Mixed: regular text alongside {\hand handwritten words} in the same line.
\end{itemize}

\subsection{Pencil Underlines}
%--------------------------------------------------

\lapis{Default underline} uses the \texttt{carandache} color (goldenrod). The color is passed as an optional argument:

\begin{democode}
\lapis{default underline}
\lapis[tomato]{custom color}
\end{democode}

\noindent The optional argument accepts any defined color (code $\to$ result):

\begin{itemize}
    \item \verb|\lapis{carandache gold}| $\to$ \lapis{carandache gold} (default)
    \item \verb|\lapis[tomato]{tomato red}| $\to$ \lapis[tomato]{tomato red}
    \item \verb|\lapis[note]{note blue}| $\to$ \lapis[note]{note blue}
    \item \verb|\lapis[cobalt]{cobalt teal}| $\to$ \lapis[cobalt]{cobalt teal}
    \item \verb|\lapis[BoEpurple]{BoE purple}| $\to$ \lapis[BoEpurple]{BoE purple}
\end{itemize}

Underlines work inside longer sentences too: the model predicts that \lapis{output falls} whenever the borrowing constraint \lapis[tomato]{binds}.

%==================================================
\section{Pencil Boxes}
%==================================================

\subsection{Inline Boxes}
%--------------------------------------------------

\lapisbox{Default gold box} uses no options. The color is set with the \texttt{color=} key:

\begin{democode}
\lapisbox{default gold box}
\lapisbox[color=tomato]{custom color}
\end{democode}

\noindent Color variants (code $\to$ result):

\begin{itemize}
    \item \verb|\lapisbox{gold}| $\to$ \lapisbox{gold} (default)
    \item \verb|\lapisbox[color=tomato]{tomato}| $\to$ \lapisbox[color=tomato]{tomato}
    \item \verb|\lapisbox[color=note]{note blue}| $\to$ \lapisbox[color=note]{note blue}
    \item \verb|\lapisbox[color=cobalt]{cobalt}| $\to$ \lapisbox[color=cobalt]{cobalt}
    \item \verb|\lapisbox[color=BoEpurple]{BoE purple}| $\to$ \lapisbox[color=BoEpurple]{BoE purple}
\end{itemize}

Boxes work mid-sentence: the key result is that \lapisbox[color=note]{$\beta > 0$} whenever agents are impatient.

\subsection{Equation Boxes}
%--------------------------------------------------

\verb|\lapisboxeq| wraps display math in a pencil box without requiring the caller to enter math mode. Put the math content directly in the argument (no \verb|$| or \verb|\[|):

\begin{democode}
\[
    \lapisboxeq{Y_t = C_t + I_t + G_t + NX_t}
    \qquad
    \lapisboxeq[color=tomato]{r_t = \rho + \phi_\pi \pi_t}
\]
\end{democode}

\noindent renders as:
\[
    \lapisboxeq{Y_t = C_t + I_t + G_t + NX_t}
    \qquad
    \lapisboxeq[color=tomato]{r_t = \rho + \phi_\pi(\pi_t - \bar\pi) + \phi_y \tilde y_t}
\]

It also works inline: \verb|\lapisboxeq{r = \rho}| $\to$ the steady-state condition is \lapisboxeq{r = \rho}.

\subsection{Labeled boxes and pencil arrows}
%--------------------------------------------------

Boxes with \verb|label=| create named TikZ nodes. A later \texttt{tikzpicture} (with the \texttt{tpstyle} option) targets them by name; compile \textbf{twice} for cross-picture references to resolve:

\begin{democode}
The \lapisbox[color=tomato,label=firesale]{firesale} externality
\begin{tikzpicture}[tpstyle]
    \draw[pencil,->,tomato!80] ([yshift=-2pt]firesale.south east)
        to[bend right] +(+0.7,-0.5)
        node[anchor=west,opacity=1,tomato] {\hand A comment};
\end{tikzpicture}
\end{democode}

\noindent renders as:

\begin{itemize}
    \item The \lapisbox[color=note,label=demand]{aggregate demand} externality
    \item The \lapisbox[color=tomato,label=firesale]{firesale} externality
\end{itemize}
\begin{tikzpicture}[tpstyle]
    \draw[pencil,->,tomato!80] ([yshift=-2pt]firesale.south east) to[bend right] +(+0.7,-0.5) node[anchor=west,opacity=1,tomato] {\hand A comment};
\end{tikzpicture}

%==================================================
\section{Annotations and Notes}
%==================================================

\subsection{Inline Todo Notes}
%--------------------------------------------------

\begin{democode}
\NB{remember to add the proof here}
\end{democode}

\noindent renders as:

\NB{This is an inline note produced by the \texttt{NB} command. It uses the \texttt{note} blue background from \texttt{todonotes}.}

Todo notes are useful for flagging gaps during drafting. They sit in the text flow (not the margin) and are easy to search.

\subsection{Draft Annotations}
%--------------------------------------------------

The default user and color are set once in the preamble; individual calls override them with \texttt{user=} and \texttt{color=} keys:

\begin{democode}
% in the preamble:
\renewcommand{\memouser}{ACB}
\renewcommand{\memocolor}{BoEred}
% in the body:
\memo{uses the preamble defaults}
\memo[user=JB,color=note]{override both}
\memo[color=cobalt]{override color only}
\end{democode}

\noindent renders as:

\memo{Default user (ACB, set once in the preamble) and default color (BoEred).}

\memo[user=JB,color=note]{Custom user and color via optional key=value argument.}

\memo[color=cobalt]{Color-only override — username still comes from \texttt{\string\memouser}.}

Regular text continues after the annotation without extra spacing.

\subsection{Note Boxes}
%--------------------------------------------------

A \verb|notebox| is a breakable shaded callout. It is an environment, so wrap content between \verb|\begin{notebox}| and \verb|\end{notebox}|:

\begin{democode}
\begin{notebox}
    \textbf{Key result.} Some highlighted text.
\end{notebox}
\end{democode}

\noindent renders as:

\begin{notebox}
    \textbf{Key result.} In the Korinek--Simsek (2016) model, competitive borrowers do not internalise that their deleveraging reduces aggregate income. The social planner internalises this and chooses a lower debt level in period 0.
\end{notebox}

Note boxes accept arbitrary content including math and itemize:

\begin{notebox}
    \textbf{Two externalities.}
    \begin{itemize}
        \item \textbf{Aggregate demand}: deleveraging $\Rightarrow$ lower income $\Rightarrow$ tighter constraint.
        \item \textbf{Firesale}: forced selling $\Rightarrow$ lower asset prices $\Rightarrow$ lower collateral.
    \end{itemize}
    Both generate over-borrowing in equilibrium.
\end{notebox}

Adding \verb|\boxentry{title}| as the first line turns a plain \verb|notebox| into a \emph{numbered} box: it gains a dotted ``Box~X.Y'' header, an entry in the List of Boxes (shown after the table of contents), and a label you can \verb|\ref|. A bare \verb|notebox| (above) stays unnumbered and header-less.

\begin{democode}
\begin{notebox}
    \boxentry{The running example}\label{box:demo}
    Body text of the numbered box ...
\end{notebox}
As shown in Box~\ref{box:demo} ...
\end{democode}

\noindent renders as:

\begin{notebox}
    \boxentry{The running example}\label{box:example}
    \textbf{The running example.} A bivariate VAR(1) keeps every object small enough to write out in full. Box numbers run ``\,arabic-section.alpha\,'' (here \ref{box:example}) and reset each section.
\end{notebox}

\begin{notebox}
    \boxentry{A second numbered box}\label{box:second}
    This box increments the letter within the section, so it is Box~\ref{box:second}. Cross-references resolve because \verb|\boxentry| calls \verb|\refstepcounter|.
\end{notebox}

Reference them anywhere: the example is in Box~\ref{box:example} and the follow-up in Box~\ref{box:second}.

Although a \verb|notebox| and an \verb|\NB| note share the same \texttt{note} blue, they serve opposite purposes. A \verb|notebox| is an environment for \emph{permanent} content you want to keep in the finished document --- a highlighted result, definition, or remark --- and it lightly shades (\texttt{note!10}), breaks across pages, and holds a body of any length. \verb|\NB{...}| is a single-argument command for \emph{transient} draft markers (``fix this later''); it renders as a more saturated inline banner (\texttt{note!40}) via \texttt{todonotes}. Unlike \verb|\memo|, \verb|\NB| has no suppression toggle, so remove such markers before finalising.

%==================================================
\section{Code Listings}
%==================================================

\subsection{The \texttt{code} Environment}
%--------------------------------------------------

The \verb|code| environment typesets a verbatim block on a light-gray background. It is deliberately monochrome --- all text is black, with no syntax highlighting --- and breaks across pages. Content is taken verbatim, so no escaping is needed:

\begin{democode}
\begin{code}
% MATLAB: estimate a reduced-form VAR(1)
[VAR, VARopt] = VARmodel(data, 1);
disp(VAR.Fcomp);   % companion matrix
\end{code}
\end{democode}

\noindent renders as:

\begin{code}
% MATLAB: estimate a reduced-form VAR(1)
root = fileparts(fileparts(mfilename('fullpath')));
addpath(fullfile(root, 'VAR'));
[VAR, VARopt] = VARmodel(data, 1);   % equation-by-equation OLS
disp(VAR.Fcomp);                     % companion matrix
rmpath(genpath(root));
\end{code}

\subsection{Inline Code Chips}
%--------------------------------------------------

\verb|\mc| and \verb|\mci| set inline code on the same light-gray background. \verb|\mc| uses a fixed code size; \verb|\mci| follows the surrounding font size:

\begin{democode}
Call \mc{addpath(genpath(root))} to add the toolbox; the
\mci{root} variable holds the install path.
\end{democode}

\noindent renders as: call \mc{addpath(genpath(root))} to add the toolbox; the \mci{root} variable holds the install path.

%==================================================
\section{Mathematics}
%==================================================

\subsection{Standard Environments}
%--------------------------------------------------

The \verb|\E| shorthand gives the expectations operator (use \verb|\E_t| for a time subscript):

\begin{democode}
\[ \E_t\left[ u'(C_{t+1}) \right] = 1 \]
\end{democode}

\noindent renders as:
\[
    \E_t\left[\frac{u'(C_{t+1})}{u'(C_t)}\frac{1+r_{t+1}}{1+\pi_{t+1}}\right] = 1
\]

A multi-line derivation using \texttt{align}:
\begin{align}
    C_t^b + B_{t+1} &= y_t + (1+r_t)B_t \\
    B_{t+1}         &\leq \phi \label{eq:constraint}\\
    \lambda_t       &= u'(C_t^b) - \beta^b(1+r_t)\E_t[u'(C_{t+1}^b)] \geq 0
\end{align}

%==================================================
\section{Tables}
%==================================================

\subsection{Standard Tabularx}
%--------------------------------------------------

The \texttt{Y} column type is a centered auto-width \texttt{X} column. The style deliberately redefines \verb|\notesize| to produce compact 8.5pt annotation text below the table:

A publication-style pattern wraps the table in a fixed-width \texttt{minipage} (so the caption and note share the table's measure), uses \verb|\shortstack| for two-line headers, and sets a long \verb|\notesize| note that wraps across several lines:

\begin{democode}
\begin{table}[ht!]
    \centering
    \begin{minipage}[b]{.7\textwidth}
        \begin{center}\small
        \begin{tabularx}{\textwidth}{lYY}
        \toprule
            & \shortstack{Demand\\ Shock} & \shortstack{Monetary\\ Policy Shock} \\
            & ($\varepsilon_{t}^{Demand}$) & ($\varepsilon_{t}^{MonPol}$) \\ \midrule\addlinespace
            Output growth ($y_{t}$)             & + & - \\
            Short-term Interest Rate ($r_{t}$)  & + & + \\ \bottomrule
        \end{tabularx}%
        \end{center}
        \vspace{0.4cm}
        \caption{\scshape{Sign Restrictions for Demand Shock \\ and Monetary Policy Shock}}
        \vspace{0.2cm}
        {\notesize {\scshape Note.} The signs represent restrictions on the
        elements of the structural impact matrix $B$ ...}
        \label{tab:SR}
    \end{minipage}
\end{table}
\end{democode}

\noindent renders as:

\begin{table}[ht!]
	\centering
	\caption{\scshape{Sign Restrictions for Demand Shock and Monetary Policy Shock}}
    \begin{minipage}[b]{.8\textwidth}
    \begin{center}\small
    \begin{tabularx}{\textwidth}{lYY}
    \toprule
        & \shortstack{Demand\\ Shock} & \shortstack{Monetary\\ Policy Shock} \\
        & ($\varepsilon_{t}^{Demand}$) & ($\varepsilon_{t}^{MonPol}$) \\ \midrule\addlinespace
        Output growth ($y_{t}$) & + & - \\
        Short-term Interest Rate ($r_{t}$) & + & + \\ \bottomrule
    \end{tabularx}%
    \end{center}
    \vspace{-.3cm}
    {\notesize {\scshape Note.} The signs represent restrictions on the elements of the structural impact matrix $B$. `$+$' and `$-$' indicate a positive and negative contemporaneous response, respectively. A demand shock raises both output growth and the short-term interest rate, as monetary policy tightens to contain the expansion. A contractionary monetary policy shock lowers output growth and raises the interest rate, consistent with an unexpected tightening of monetary policy.\par}
    \label{tab:SR}
	\end{minipage}
\end{table}
 
\subsection{Table with Row Colors}
%--------------------------------------------------

The \texttt{[table]} option on \texttt{xcolor} (loaded by the theme) enables \verb|\rowcolor|, placed at the start of a row:

\begin{democode}
\rowcolor{note!8} Discount factor & $\beta^b$ & 0.90 \\
\end{democode}

\noindent renders as:

\begin{table}[ht!]
	\centering
	\caption{\scshape{Model Parameters}}
    \begin{minipage}[b]{.7\textwidth}
    \begin{center}\small
    \begin{tabularx}{\textwidth}{lYY}
    \toprule
    \textbf{Parameter} & \textbf{Symbol} & \textbf{Value} \\
    \midrule
    \rowcolor{note!8} Discount factor (borrower) & $\beta^b$ & 0.90 \\
    Discount factor (lender)  & $\beta^l$ & 0.95 \\
    \rowcolor{note!8} Borrowing limit           & $\phi$    & 0.75 \\
    Frisch elasticity         & $\varphi$ & 1.00 \\
    \rowcolor{note!8} Price rigidity            & $\kappa$  & 0.50 \\
    \bottomrule
    \end{tabularx}%
    \end{center}
    \vspace{-.3cm}
    {\notesize {\scshape Note.} Illustrative calibration of the borrower--lender model used throughout this demo; the values are not estimated and serve only to exercise the \texttt{\string\rowcolor} shading and the \texttt{Y} column type within the same fixed-width \texttt{minipage} layout as Table~\ref{tab:SR}.\par}
    \label{tab:params}
	\end{minipage}
\end{table}

%==================================================
\section{Paragraph Numbering}
%==================================================

Numbered paragraphs use \verb|\para|. Place it at the start of the paragraph; the counter increments automatically and typesets a number followed by a quad of space:

\begin{democode}
\para This paragraph is automatically numbered.
\para So is this one, with the next number.
\end{democode}

\noindent renders as:

\para This is the first numbered paragraph. The number is produced by \verb|\para| and counts independently of section numbers. It is useful for notes that cross-reference specific arguments.

\para\label{para:bind} This is the second numbered paragraph. The borrowing constraint $B_{t+1} \leq \phi$ is assumed to bind in period 1 but not in period 0.

\para This is the third numbered paragraph. The social planner internalises the aggregate demand externality and imposes a tax $\tau^*$ on period-0 borrowing.

Because \verb|\para| calls \verb|\refstepcounter|, a numbered paragraph can carry a \verb|\label| and be cited later with \verb|\ref|:

\begin{democode}
\para\label{para:key} The constraint binds in period 1.
...as assumed in paragraph~\ref{para:key}, the planner...
\end{democode}

\noindent The second paragraph above is labelled \verb|para:bind|; referring to it here resolves to paragraph~\ref{para:bind}. A fourth paragraph can then build on it explicitly:

\para Building on the assumption introduced in paragraph~\ref{para:bind}, the planner's tax $\tau^*$ restores constrained efficiency.

\begin{notebox}
    Cross-references resolve inside environments too: this \verb|notebox| cites paragraph~\ref{para:bind} directly, confirming that \verb|\para| labels behave like any other counter and that \verb|\para| coexists with \verb|notebox| without disturbing the count.
\end{notebox}

%==================================================
\section{Pencil Arrows}
%==================================================

\subsection{Inline Directional Arrows}
%--------------------------------------------------

The \verb|\jarrow*| commands produce small baseline-aligned pencil arrows. Follow them with \verb|\ | (backslash-space) to keep the following space, since the command gobbles it:

\begin{democode}
GDP growth \jarrowup\ accelerated in Q3.
\jarrowup[color=tomato]   % options via key=value
\jarrowr[len=0.6cm,lw=0.1cm]   % \jarrowr = \jarrowright
\end{democode}

\noindent renders as (all four directions are available):

\begin{itemize}
    \item Up:    GDP growth \jarrowup\ accelerated in Q3.
    \item Down:  Inflation \jarrowdown\ fell below target.
    \item Right: The effect propagates \jarrowright\ through the banking sector.
    \item Left:  Demand shifted \jarrowleft\ after the shock.
\end{itemize}

Custom options via key=value: \jarrowup[color=tomato]\ (tomato), \jarrowr[len=0.6cm,lw=0.1cm]\ (long, thick), \jarrowd[color=note,pad=0.1em]\ (tight blue).

\subsection{Overlay Arrows with Marker Nodes}
%--------------------------------------------------

\verb|\marker{name}{text}| wraps visible text in a named TikZ node. A subsequent overlay \texttt{tikzpicture} draws an \texttt{arrow}-style line to it. Compile \textbf{twice} for cross-picture references to resolve:

\begin{democode}
Consider the \marker{fisherA}{Fisher equation}:
\begin{tikzpicture}[tpstyle]
    \draw[arrow] ([yshift=-3pt]fisherA.south) to[bend right=20]
        +(+0.2,-0.6) node[anchor=north west] {\hand links $r$ and $\pi$};
\end{tikzpicture}
\end{democode}

\noindent renders as:

Consider the \marker{fisherA}{Fisher equation} and the \marker{quantA}{quantity theory}:
\begin{tikzpicture}[tpstyle]
    \draw[arrow] ([yshift=-3pt]fisherA.south) to[bend right=20] +(+0.2,-0.6)
        node[anchor=north west] {\hand links $r$ and $\pi$};
    \draw[arrow,note] ([yshift=-3pt]quantA.south) to[bend right=20] +(-0.2,-0.6)
        node[anchor=north west] {\hand links $M$ and $P$};
\end{tikzpicture}
\vspace{1.5cm}

%==================================================
\section{Package Options}
%==================================================

This section documents the two options the package accepts at load time, passed in the optional argument of \verb|\usepackage|: \verb|wiggle=<n>| sets the amplitude of the pencil decoration, and \verb|hidememo| suppresses all \verb|\memo| output for a clean final build. Each is demonstrated below.

\subsection{Wiggle}
%--------------------------------------------------

The pencil amplitude is controlled by the \verb|wiggle=| package option:

\begin{itemize}
    \item \verb|\usepackage{jambronotes}| $\to$ default amplitude (0.4pt)
    \item \verb|\usepackage[wiggle=0.3]{jambronotes}| $\to$ subtle
    \item \verb|\usepackage[wiggle=1.5]{jambronotes}| $\to$ expressive
\end{itemize}

This document uses the default. Recompile with a different option to compare:
\[
    \lapisboxeq{Y = \alpha + \beta X + \varepsilon} \qquad
    \lapis{underline at default wiggle}
\]

\subsection{Memo Suppression}
%--------------------------------------------------

Use \verb|\usepackage[hidememo]{jambronotes}| when you want the same notes document without draft annotations. In that mode every \verb|\memo| call expands to nothing, while the rest of the style remains unchanged.

%==================================================
\section{Front and Back Matter}
%==================================================

\subsection{TOC Headings and the List of Boxes}
%--------------------------------------------------

The List of Boxes after the table of contents is produced by \verb|\listofboxes|; every numbered box (created with \verb|\boxentry|, Section~III) adds an entry. \verb|\tocheading| typesets an unnumbered heading in the TOC font, useful for labelling front-matter sections:

\begin{democode}
\tableofcontents
\listofboxes              % auto-titled "List of Boxes"
\tocheading{Supplementary Material}   % unnumbered heading
\end{democode}

\noindent \verb|\tocheading{Example Heading}| renders as:

\tocheading{Example Heading}

\subsection{Appendix Numbering}
%--------------------------------------------------

\verb|\backmatter| switches \verb|\section| numbering to letters and restarts it at~A (and re-prefixes hyperlink anchors so they stay unique), matching the appendix convention of the VAR Toolbox Handbook. Call it before an appendix; the section below uses it, so it is lettered~A even though it follows Section~X.

\begin{democode}
\backmatter
\section{Appendix: Quick Reference}
\end{democode}

\backmatter
%==================================================
\section{Appendix: Quick Reference}
%==================================================

This section follows Section~X in the source but is lettered~\thesection, because \verb|\backmatter| switched section numbering to letters and reset the counter. 

\end{document}
